\subsection*{Integers for Program Control}

In high-performance computation kernels, integers are typically used for program control rather than computation. In Croqtile \emph{tileflow programs}, integers are used exclusively for loop control and array indexing. Given the limited value ranges in these scenarios, Croqtile has provided a single integer type --- the \emph{32-bit signed integer} --- to simplify the design for this domain-specific use.

In Croqtile, an \textbf{integer}-typed variable is defined similar to C/C++:

\begin{lstlisting}[language=croqtile,caption={Integers for Program Control}]
int a = 1;
\end{lstlisting}

However, since Croqtile can infer the type from the \emph{initialization expression}, you can also define an integer variable in another style:

\begin{lstlisting}[language=croqtile,caption={Integers for Program Control}]
a = 1;
\end{lstlisting}

Similar to C/C++, you can define multiple integers in a single line:

\begin{lstlisting}[language=croqtile,caption={Integers for Program Control}]
int a = 1, b = 2, c = 3;
d = 4, e = 5;
\end{lstlisting}

In Croqtile, the \emph{mdspan}-typed shape can not be used as the function parameter. However, integers can be passed from the host to the Tileflow program, and from the Tileflow program to the device:

\begin{lstlisting}[language=croqtile,caption={Integers for Program Control}]
__co__ void foo(int a) {
  // ...
  call kernel(a, 3);
}
\end{lstlisting}

Unlike C/C++, Croqtile does not allow the declaration of integers without initialization or the reassignment of an integer. For example:

\begin{lstlisting}[language=croqtile,caption={Integers for Program Control}]
int a;  // error: declaration without initialization
int b = 1;
b = 3;  // error: re-assignment
\end{lstlisting}

This enforces that \textbf{integers must be initialized at declaration and cannot be reassigned}. In some contexts, these integers are referred to as "\textbf{immutable}" or "\textbf{constant}".
